\section{The Propers: Themes and Movement}
\zlabel{triptych:concise:themes:start}
The Lord leaves time for people to repent and teaches them to ask for the
salvation he promises. The Gospel explains why the servants must spare
people they would uproot; the prayers concern people who themselves need
protection, pardon and the completion of a gift already received.

\subsection{Peace promised during exile}
The Introit begins with God's assurance of peace to the captives in Babylon
(Jer 29:11,12,14). In the letter, they must build houses, plant and raise
families while awaiting return. Jerome opposes the false prophets' quick
return to God's command. He then explains why the promise still calls for
prayer: God urges the exiles to ask so that they may receive the good he
could give without their invocation. The invitation concerns their
obedience as recipients.\footnote{Jerome, \work{Commentary on Jeremias},
V, on 29:10--14, PL 24, cols.~893--894.}

The psalm verse thanks God for restoring Jacob (Ps 84:2), although the
whole psalm continues to ask for restoration. Augustine explains this
combination through a promised grace that the prophet sees and still
asks to receive. In his Christian exposition, captivity includes bondage
under sin, and release includes pardon.\footnote{Augustine,
\work{Exposition of Psalm 84}, §§2--6, NPNF heading Psalm 85.}
The Collect asks continuing protection for God's family. Read beside
these texts, its reliance on heavenly grace describes dependence that
continues after God has begun to give.

\subsection{The forgiven community and its grievances}
Paul addresses chosen, holy and beloved people who still have complaints
against each other (Col 3:12--17). He commands forgiveness by pointing to
the Lord's forgiveness of them. Chrysostom explains that even a serious
injury falls within that example. He calls charity the bond without which
apparently good conduct comes apart, as a house needs fastenings and bones
need ligaments. His image of peace as umpire describes a grievance in
which revenge and endurance compete; the prize belongs to endurance.
The same Epistle commands teaching and admonition, giving the community
work to do while it bears with the offender.\footnote{Chrysostom,
\work{Homilies on Colossians}, VIII, on 3:12--15; IX, opening on 3:16--17.
His lemma has ``peace of God'' where the Missal has ``peace of Christ.''}

Before the Gospel, the Gradual thanks God for rescue (Ps 43:8--9), while
the Alleluia asks him to hear the cry from the depths (Ps 129:1--2).
Psalm 43 also laments defeat. Augustine hears the martyrs in it and
explains how assured future rescue can be spoken of as accomplished.
Praise can therefore accompany present affliction. In Psalm 129 he
recognises hope in the sinner who cries: the petitioner knows his need
and expects mercy.\footnote{Augustine, \work{Exposition of Psalm 43},
§§1,3--8; \work{Exposition of Psalm 129}, §§1--2.}

Paul's commands also concern the life from which these people have turned.
Colossians 3 opposes anger, malice and lying to the mercy they must put on.
The new life requires changed conduct from those already called holy and
beloved. Chrysostom explains the abundance of Christ's word as instruction
available in ordinary households, including times of loss. Its consolation
continues after the immediate trouble. The assembly approaches the Gospel
as people being taught, with grievances to forgive and songs of thanks
to offer. Their own need for instruction will matter when the servants
propose to correct the field.

\clearpage
\subsection{The crop that hasty zeal could destroy}
The Gospel's servants correctly recognise that an enemy has sown tares.
Their next proposal raises another question: may they pull them out?
The owner refuses because wheat could be uprooted with them
(Matt 13:24--30). In Jesus's later explanation, the seed signifies people,
the field is the world and the final reapers are angels (13:36--43).
The owner's reason concerns the people endangered by the proposed remedy.

Chrysostom identifies the tares with heretics and explains the danger in
two ways. War against them can kill saints; some people presently in error
may also become wheat through repentance. He forbids killing in this
argument while allowing nonlethal suppression of speech and assemblies.
Augustine addresses wicked Christians more broadly. People may change;
human beings may also mistake wheat for tares, whereas the angels do not
err. He therefore warns hearers who think themselves good as well as
those who need to repent.\footnote{Chrysostom, \work{Homilies on Matthew},
XLVI.1; Augustine, Sermon XXIII [Benedictine LXXIII], §§1--4.}

These explanations concern restraint toward a neighbour whose future
remains open. Paul's command to forgive and admonish describes active
care during that time. His words also address the person who prays:
forgiveness received becomes forgiveness given. The two interpretations
thus meet in the same community, while asking different questions about
its care for others and its dependence on God.

\subsection{From the depths to a pledge of salvation}
At the Offertory the cry from Psalm 129 returns. The Secret then asks
absolution for the offerers' offences and guidance for their wavering
hearts. Augustine explains the psalm's propitiation through Christ's
sacrifice for sinners. Connecting that exposition with the Secret,
the offerers can ask pardon because Christ has acted for the guilty.
Their request for direction concerns the life that follows pardon.
The Preface gives thanks to the Father with the Son and Holy Spirit,
confessing distinct Persons, one essence and equal majesty. Thanksgiving
through Christ, commanded in Colossians, is offered to this God.

The Communion appoints Christ's assurance to believing prayer
(Mark 11:24). Bede explains it beside the following command to forgive:
faith works through charity, and those who need pardon can ask for it
while forgiving their offenders. Aquinas distinguishes the merit of
prayer from God's gift of what it obtains; a faithfully requested temporal
good may be withheld for the petitioner's salvation.\footnote{Bede,
\work{Exposition of Mark}, on 11:22--26; Aquinas, \work{Summa theologiae}
II--II, q.~83, aa.~15--16. Aquinas supplies doctrinal exposition of prayer.}
The Postcommunion specifies a saving request. People who have received
salvation's pledge ask to attain its effect. The harvest promises final
gathering and judgement; this prayer asks that the present recipients
reach salvation's completion.

Augustine explains the waiting of Psalm 129 through Christ's resurrection.
The flesh he took from humanity has risen, giving his members reason to
hope for their own resurrection.\footnote{Augustine, \work{Exposition of
Psalm 129}, §§3--8.} In relation to the Postcommunion, this is a hope
for completed salvation grounded in a gift already made. In relation to
the Gospel, the remaining time permits repentance before final separation.
The first interpretation consequently asks how to act toward a person
who can still change. The second asks how that same needy person can
approach God and persevere in expectation of the promised end.
\zlabel{triptych:concise:themes:end}
